\subsubsection{Commutator norm of the arithmetic projections}

The electronic prefactor is obtained without using \(\alpha\), torsion,
or a dimensional unit. The domain is the third realization of APP,
\(\mathcal A_3=\{1,\ldots,9\}^3\), equipped with the uniform measure.
The function space is \(\ell^2(\mathcal A_3;\mathbb R)\), with inner
product
\[
 \langle f,g\rangle=\frac1{729}\sum_{x\in\mathcal A_3}f(x)g(x).
\]
In this block, \([P,Q]=PQ-QP\) denotes the additive operator commutator.
We define
\[
 \Sigma(x)=\rho_9(x_1+x_2+x_3),\qquad
 \Pi(x)=\rho_9(x_1x_2x_3).
\]
The conditional expectations \(P_\Sigma\) and \(P_\Pi\) are the
orthogonal projections onto functions constant on each additive
and multiplicative fiber, respectively. Thus, for a function \(f\),
\[
 (P_\Sigma f)(x)
 =\frac{1}{|\Sigma^{-1}(\Sigma(x))|}
   \sum_{y:\,\Sigma(y)=\Sigma(x)}f(y),
\]
and analogously for \(P_\Pi\). The two projections act on the
same space; they do not represent two independent arithmetic supports.

\begin{lemma}[Two projections and principal angles]
If \(s\in[0,1]\) is a singular value of the pairing between
orthonormal bases of the ranges of two projections \(P,Q\), its nontrivial
block contributes \(s\sqrt{1-s^2}\) to
\(\|[P,Q]\|\).
\end{lemma}
\begin{proof}
On the corresponding principal plane, one can write
\[
 P=\begin{pmatrix}1&0\\0&0\end{pmatrix},\qquad
 Q=\begin{pmatrix}
 s^2&s\sqrt{1-s^2}\\
 s\sqrt{1-s^2}&1-s^2
 \end{pmatrix}.
\]
Their commutator is
\[
 [P,Q]=s\sqrt{1-s^2}
 \begin{pmatrix}0&1\\-1&0\end{pmatrix}.
\]
The final matrix is orthogonal, so its norm is one. Common
or orthogonal blocks have zero commutator. The orthogonal decomposition
into principal planes gives the claim and reduces the total norm to the maximum
of these contributions.
\end{proof}

\begin{proposition}[Exact incompatibility norm]
On \(\ell^2(\mathcal A_3)\),
\begin{equation}
 \|[P_\Sigma,P_\Pi]\|=\frac{\sqrt3}{4}.
 \label{ivloc:norma}
\end{equation}
\end{proposition}
\begin{proof}
Let \(N_{ab}=|\Sigma^{-1}(a)\cap\Pi^{-1}(b)|\). Integer enumeration
of the three symbols yields
\[
 N=9\begin{pmatrix}
 0&1&1&0&1&1&0&1&4\\
 1&0&1&1&0&1&1&0&4\\
 1&0&2&0&1&2&0&0&3\\
 0&1&1&0&1&1&0&1&4\\
 1&0&1&1&0&1&1&0&4\\
 0&0&2&1&0&2&0&1&3\\
 0&1&1&0&1&1&0&1&4\\
 1&0&1&1&0&1&1&0&4\\
 0&1&2&0&0&2&1&0&3
 \end{pmatrix}.
\]
Each row sums to \(81\); the column sums are
\[
 (m_1,\ldots,m_9)=(36,36,108,36,36,108,36,36,297).
\]
The normalized indicator functions of the fibers have pairing
matrix \(C_{ab}=N_{ab}/\sqrt{81m_b}\). Consequently,
\[
 M=CC^{\mathsf T},\qquad
 M_{ac}=\sum_{b=1}^9\frac{N_{ab}N_{cb}}{81m_b}
\]
is a rational matrix whose spectrum consists of the squares of the
required singular values.

Exact multiplication shows that \((1,\ldots,1)\),
\((1,-1,0)^3\), and \((1,1,-2)^3\) are eigenvectors with eigenvalues
\(1\), \(1/4\), and \(7/132\), respectively. The subspace of
vectors supported on \(\{3,6,9\}\) whose sum is zero has dimension
two and eigenvalue \(1/18\). The zero-sum vectors supported on
\(\{1,4,7\}\) or \(\{2,5,8\}\) form a four-dimensional
kernel. These subspaces are independent and their dimensions sum to
nine. The complete spectrum is therefore established:
\[
 \operatorname{spec}(M)=
 \left\{1,\frac14,\frac7{132},\frac1{18},\frac1{18},0,0,0,0\right\}.
\]
The maximum of \(\lambda(1-\lambda)\) on this set is \(3/16\),
attained at \(\lambda=1/4\). The preceding lemma gives
\(\sqrt{3/16}=\sqrt3/4\).
\end{proof}

The mode \((1,-1,0)^3\) is precisely the distribution of values of the
tritic phase selector in the nonadic chart. The scalar norm retains
a measure of incompatibility, but does not by itself preserve the
orientation of that mode or reconstruct the two projectors. That memory
remains in the route state accompanying the scalar reading.

\begin{proposition}[Algebra of the local electronic block]
\label{ivloc:algebra}
On the principal plane attaining \eqref{ivloc:norma}, let
\[
 P=\begin{pmatrix}1&0\\0&0\end{pmatrix},\qquad
 Q=\begin{pmatrix}1/4&\sqrt3/4\\\sqrt3/4&3/4\end{pmatrix},\qquad
 S=2P-I,\qquad J=\frac4{\sqrt3}[P,Q].
\]
Then \(S^2=I\), \(J^2=-I\), \(SJ=-JS\), and the real algebra
generated by \(S,J\) is \(M_2(\mathbb R)\), a realization of
\(\mathrm{Cl}_{1,1}\). Moreover,
\[
 n_\pm=\frac{S\pm J}{2},\qquad
 n_\pm^2=0,\qquad n_+n_-+n_-n_+=I.
\]
\end{proposition}
\begin{proof}
We have \(S=\operatorname{diag}(1,-1)\) and
\(J=\bigl(\begin{smallmatrix}0&1\\-1&0\end{smallmatrix}\bigr)\).
Multiplication verifies the three relations. The matrices
\(I,S,J,SJ\) are linearly independent and form a basis of
\(M_2(\mathbb R)\), proving the claim about the algebra.
Finally, \((S\pm J)^2=S^2+J^2\pm(SJ+JS)=0\), and the sum of the
cross products is \((S^2-J^2)/2=I\).
\end{proof}

In this block, an elliptic direction, a hyperbolic direction, and two
parabolic null directions coexist. The result describes a local
matrix representation of the arithmetic of the two sheets. It does not identify all of APP
with \(M_2(\mathbb R)\), or with a quantum Latin square; such
claims would require other domains and other relations. Nor does the
exponential belong exclusively to the parabolic regime: it can
be applied to any of the generators of this block.
